% chktex-file 44
% chktex-file 3
\documentclass{ctexart}
\usepackage[top=2cm, bottom=2cm, left=2.5cm, right=2.5cm]{geometry}
\usepackage{graphicx}
\graphicspath{{D:/code/LATEX/Figures}}
\usepackage{amsmath}
\usepackage{amssymb} 
\usepackage{hyperref}
\usepackage{fix-cm}
\usepackage{bm}
\usepackage{booktabs}
\usepackage{array}
\vfuzz=100pt  % 垂直方向容忍度
\hfuzz=100pt  % 水平方向容忍度  
\vbadness=10000
\hbadness=10000
\overfullrule=0pt  % 不标记过满行

\begin{document}
\section{Dini导数}
函数$f(x)$的四个迪尼导数如下：
\begin{align}
    D^+f(x)=\varlimsup_{h\to 0^+}\frac{f(x+h)-f(x)}{h} \\
    D_+f(x)=\varliminf_{h\to 0^+}\frac{f(x+h)-f(x)}{h} \\
    D^-f(x)=\varlimsup_{h\to 0^-}\frac{f(x+h)-f(x)}{h} \\
    D_-f(x)=\varliminf_{h\to 0^-}\frac{f(x+h)-f(x)}{h}
\end{align}
它们依次称为右上导数、右下导数、左上导数、左下导数。显然
\[D_+f(x)\leq D^+f(x),D_-f(x)\leq D^-f(x)\]

（命题1.1）迪尼导数不变号能够推出单调性：
\begin{enumerate}
    \item 设$f(x)\in C[0,1]$，且任意$x\in (0,1)$有$D^+f(x)>0$，则$f(x)$严格单调递增；\\
          设$f(x)\in C[0,1]$，且任意$x\in (0,1)$有$D^+f(x)\geq 0$，则$f(x)$单调递增
    \item 设$f(x)\in C[0,1]$，且任意$x\in (0,1)$有$D_+f(x)>0$，则$f(x)$严格单调递增；\\
          设$f(x)\in C[0,1]$，且任意$x\in (0,1)$有$D_+f(x)\geq 0$，则$f(x)$单调递增
    \item 设$f(x)\in C[0,1]$，且任意$x\in (0,1)$有$D^-f(x)<0$，则$f(x)$严格单调递减；\\
          设$f(x)\in C[0,1]$，且任意$x\in (0,1)$有$D^-f(x)\leq 0$，则$f(x)$单调递减
    \item 设$f(x)\in C[0,1]$，且任意$x\in (0,1)$有$D_-f(x)<0$，则$f(x)$严格单调递减；\\
          设$f(x)\in C[0,1]$，且任意$x\in (0,1)$有$D_-f(x)\leq 0$，则$f(x)$单调递减
\end{enumerate}
证明：
1.$D^+f(x)>0$时，$\forall a,b\in (0,1),a<b$,若$f(a)\geq f(b)$,则由
\[D^+f(x_0)=\varlimsup_{h\to 0^+}\frac{f(x_0+h)-f(x_0)}{h}>0\]
知$f(x)$在区间$[a,b]$上的最大值在某个内点$x_0$处取到，$D^+f(x_0)\leq 0$，矛盾。\\
$D^+f(x)\geq 0$时，$\forall a,b\in (0,1),a<b$,若$f(a)\geq f(b)$，
考虑取$\epsilon>0$充分小，使$g(x)=f(x)+\epsilon x$满足
$g(a)>g(b)$，则$D^+g(x)=D^+f(x)+\epsilon>0$，化归为上一种情况。\\
2.$D_+f(x)\leq D^+f(x)$，化归为1的证明。\\
3.考虑$g(x)=f(1-x)$,
\begin{align*}
    D^+g(x) & =\varlimsup_{h\to 0^+}\frac{g(x+h)-g(x)}{h}      \\
            & =\varlimsup_{h\to 0^+}\frac{f(1-x-h)-f(1-x)}{h}  \\
            & =-\varlimsup_{h\to 0^-}\frac{f(1-x+h)-f(1-x)}{h} \\
            & =-D^-f(1-x)>0
\end{align*}
所以$g(x)$严格单调递增，$f(x)$严格单调递减。\\
4.与3同理。\\
只要考虑$h(x)=-f(x)$，即可证明在剩余的情况下$f(x)$也有单调性。

（命题1.2）如果迪尼导数中有一个恒为0，则$f(x)$是常值函数。\\
证明：迪尼导数（以$D^+f(x)$为例）恒为0能够推出$D^+f(x)\leq 0$和$D^+f(x)\geq 0$，
根据命题1.1，$f(x)$既单调递增又单调递减，所以是常函数。

（命题1.3）Dini导数的罗尔中值定理：设$f\in C[a,b],f(a)=f(b)$,则\[\exists\xi\in (a,b),D_+f(\xi)\leq D^+f(\xi)\leq 0\leq D_-f(\xi)\leq D^-f(\xi)\]
或\[\exists\eta\in (a,b),D^+f (\eta)\leq D_+f (\eta)\leq 0\leq D^-f(\eta)\leq D_-f(\eta)\]
证明：连续函数f在闭区间[a,b]上必有最值点，如果最大值和最小值均为$f(a)=f(b)$，则
f为常函数，迪尼导数恒为0；如果最大值不为$f(a)=f(b)$，记f在$\xi\in (a,b)$处取
到最大值，则对充分小的h恒有$f(\xi +h)\leq f(\xi)$，此时
$D_+f(\xi)\leq D^+f(\xi)\leq 0\leq D_-f(\xi)\leq D^-f(\xi)$;
如果最小值不为$f(a)=f(b)$，记f在$\eta\in (a,b)$处取到最小值，则对充分小的h恒有
$f(\eta +h)\geq f(\eta)$，此时$D^+f (\eta)\leq D_+f (\eta)\leq 0\leq D^-f(\eta)\leq D_-f(\eta)$\\

（命题1.4）Dini导数的拉格朗日中值定理：设$f\in C[a,b]$，则
\[\exists\xi\in (a,b),D_+f(\xi)\leq D^+f(\xi)\leq \frac{f(b)-f(a)}{b-a}\leq D_-f(\xi)\leq D^-f(\xi)\]
或\[\exists\eta\in (a,b),D^+f (\eta)\leq D_+f (\eta)\leq \frac{f(b)-f(a)}{b-a}\leq D^-f(\eta)\leq D_-f(\eta)\]
证明：考虑辅助函数$g(x)=f(x)-\frac{f(b)-f(a)}{b-a}x$,则
$g(a)=g(b)=\frac{b}{b-a}f(a)-\frac{a}{b-a}f(b)$，用命题1.3即证。

当x为f的第一类间断点时，仍然可以定义迪尼导数：
\begin{align}
    D^+f(x)=\varlimsup_{h\to 0^+}\frac{f(x+h)-f(x+)}{h} \\
    D_+f(x)=\varliminf_{h\to 0^+}\frac{f(x+h)-f(x+)}{h} \\
    D^-f(x)=\varlimsup_{h\to 0^-}\frac{f(x+h)-f(x-)}{h} \\
    D_-f(x)=\varliminf_{h\to 0^-}\frac{f(x+h)-f(x-)}{h}
\end{align}
其中$f(x+)=\lim_{h\to 0+}f(x+h),f(x-)=\lim_{h\to 0-}f(x+h)$

\section{傅里叶级数的逐点收敛性}
设函数$f(t)$以T为周期且$f\in L^1([0,T])$，则指数形式的傅里叶系数
\[c_n=\frac{1}{T}\int_{0}^{T}f(t)e^{-in\omega t}\,dt\]存在，并且是良定义的：
\[|c_n|\leq \frac{1}{T}\int_{0}^{T}|f(t)e^{-in\omega t}|\,dt=\frac{1}{T}\int_{0}^{T}f(t)\,dt<\infty \]

为了研究傅里叶级数$\sum_{n=-\infty}^{\infty}c_n e^{in\omega t}$在何时收敛于$f(t)$，
引入狄利克雷核\[D_N(t)=\sum_{k=-N}^{N}e^{ik\omega t}=1+\sum_{k=1}^{N}(e^{ik\omega t}+e^{-ik\omega t})=1+2\sum_{k=1}^{N}\cos(k\omega t)\]
将傅里叶级数的部分和记作\[S_N^f(t)=\sum_{-N}^{N}c_k e^{ik\omega t}\]
则有恒等式\begin{align*}
    S_N^f(t) & =\sum_{-N}^{N}c_k e^{ik\omega t}                                                  \\
             & =\sum_{-N}^{N}(\frac{1}{T}\int_{T}f(\tau)e^{-k\omega \tau}\,d\tau) e^{ik\omega t} \\
             & =\frac{1}{T}\int_{T}f(\tau)\sum_{k=-N}^{N}e^{ik\omega (t-\tau)}\,d\tau            \\
             & =\frac{1}{T}\int_{T}f(\tau)D_N(t-\tau)\,d\tau                                     \\
             & =\frac{1}{T}\int_{T}f(t-\tau)D_N(\tau)\,d\tau                                     \\
             & =\frac{1}{T}\int_{T}f(t+\tau)D_N(\tau)\,d\tau
\end{align*}
可以用等比数列求和或积化和差裂项的方法化简$D_N(t)$：
\begin{align*}
    D_N(t) & =\sum_{k=-N}^{N}e^{ik\omega t}=e^{-iN\omega t}\frac{1-e^{i(2N+1)\omega t}}{1-e^{i\omega t}}                  \\
           & =\frac{e^{i(N+1)\omega t}-e^{-iN\omega t}}{e^{i\omega t}-1}                                                  \\
    D_N(t) & =1+\sum_{k=1}^{N}(e^{ik\omega t}+e^{-ik\omega t})=1+2\sum_{k=1}^{N}\cos(k\omega t)                           \\
           & =1+\frac{2}{\sin(\frac{\omega t}{2})}\sum_{k=1}^{N}\cos(k\omega t)\sin(\frac{\omega t}{2})                   \\
           & =1+\frac{1}{\sin(\frac{\omega t}{2})}\sum_{k=1}^{N}(\sin(k+\frac{1}{2})\omega t-\sin(k-\frac{1}{2})\omega t) \\
           & =1+\frac{\sin(N+\frac{1}{2})\omega t-\sin(\frac{\omega t}{2})}{\sin(\frac{\omega t}{2})}                     \\
           & =\frac{\sin(N+\frac{1}{2})\omega t}{\sin(\frac{\omega t}{2})}
\end{align*}

（命题2.1）\begin{align}
    \int_{-\frac{T}{2}}^{0}D_N(t)\,dt=\int_{0}^{\frac{T}{2}}D_N(t)\,dt=\frac{T}{2}
\end{align}
证明：
\begin{align*}
    \int_{0}^{\frac{T}{2}}D_N(t)\,dt & =\int_{0}^{\frac{T}{2}}(1+\sum_{k=1}^{N}\cos(k\omega t))\,dt                  \\
                                     & =\frac{T}{2}+\sum_{k=1}^{N}\frac{\sin(k\omega t)}{k\omega}|_{0}^{\frac{T}{2}} \\
                                     & =\frac{T}{2}+\frac{1}{\omega}\sum_{k=1}^{N}\frac{\sin(k\pi)}{k}=\frac{T}{2}   \\
    \intertext{$D_N(t)$是偶函数，得证。}
\end{align*}

（命题2.2）Bessel不等式：设$ f\in L^2([0,T]),c_n=\frac{1}{T}\int_{T}f(t)e^{-ik\omega t}\,dt$，则
\begin{align}
    \sum_{-\infty}^{\infty}|c_n|^2\leq \frac{1}{T}\int_{T}|f(t)|^2\,dt
\end{align}
证明：
\begin{align*}
    |f(t)-\sum_{n=-N}^{N}c_n e^{in\omega t}|^2 & =(f(t)-\sum_{n=-N}^{N}c_n e^{in\omega t})(f(t)-\sum_{n=-N}^{N}c_n e^{in\omega t})^*                                       \\
                                               & =(f(t)-\sum_{n=-N}^{N}c_n e^{in\omega t})(f^*(t)-\sum_{n=-N}^{N}c_n e^{-in\omega t})                                      \\
                                               & =|f(t)|^2-\sum_{n=-N}^{N}(c_n^*f(t)e^{in\omega t}+c_n f^*(t)e^{-in\omega t})+\sum_{m,n=-N}^{N}c_m c_n^*e^{i(m-n)\omega t}
\end{align*}
将上式在一个周期上积分，由于
\[\int_{T}f(t)e^{in\omega t}\,dt=Tc_n,\int_{T}e^{i(m-n)\omega t}\,dt=\begin{cases}
        0 & \text{if }m\neq n \\
        T & \text{if }m=n
    \end{cases}\]
故\begin{align*}
      & \int_{T}|f(t)|^2\,dt-\sum_{n=-N}^{N}(c_n^*\int_{T}f(t)e^{in\omega t}\,dt+c_n \int_{T}f^*(t)e^{-in\omega t}\,dt)+\sum_{m,n=-N}^{N}c_m c_n^*\int_{T}e^{i(m-n)\omega t}\,dt \\
    = & \int_{T}|f(t)|^2\,dt-T\sum_{n=-N}^{N}(c_n^* c_n+c_n c_n^*)+T\sum_{n=-N}^{N}c_n^* c_n                                                                                     \\
    = & \int_{T}|f(t)|^2\,dt-T\sum_{n=-N}^{N}|c_n|^2
\end{align*}
这是非负函数的积分，积分值非负，即\begin{align*}
    \sum_{-\infty}^{\infty}|c_n|^2\leq \frac{1}{T}\int_{T}|f(t)|^2\,dt<\infty
\end{align*}

（定义2.3）称函数$f(t)$在[a,b]上是分段连续的并记$f\in PC[a,b]$，如果
\begin{itemize}
    \item 除了有限个点$a<x_1<x_2<\cdots<x_k<b$外，f是连续的
    \item $\forall j\in\{1,2,\dots,k\},f(x_j-)=\lim_{h\to 0^-}f(x_j+h)$和$f(x_j+)=\lim_{h\to 0^+}f(x_j+h)$存在
    \item $f(a+)=\lim_{h\to 0^+}f(a+h)$和$f(b-)=\lim_{h\to 0^-}f(b+h)$存在
\end{itemize}

（命题2.4）如果$f(t)\in L^1(\mathbb{R})$分段连续，且在点$t_0$处，四个
迪尼导数均有限，即
\begin{align*}
    \varlimsup_{t \to 0^+}\frac{f(t_0+t)-f(t_0+)}{t}<\infty,\varliminf_{t \to 0^+}\frac{f(t_0+t)-f(t_0+)}{t}<\infty \\
    \varlimsup_{t \to 0^-}\frac{f(t_0+t)-f(t_0-)}{t}<\infty,\varliminf_{t \to 0^-}\frac{f(t_0+t)-f(t_0-)}{t}<\infty \\
\end{align*}
则\[\lim_{N\to\infty}S_N^f(t_0)=\frac{f(t_0+)+f(t_0-)}{2}\]
证明：取迪尼导数的共同的上界M>0.由命题（2.1），
\begin{align*}
    S_N^f(t_0)-\frac{f(t_0^+)+f(t_0^-)}{2} & =\frac{1}{T}(\int_{T}f(t_0-\tau)D_N(\tau)\,d\tau-\int_{0}^{\frac{T}{2}}f(t_0^+)D_N(\tau)\,d\tau-\int_{-\frac{T}{2}}^{0}f(t_0^-)D_N(\tau)\,d\tau) \\
                                           & =\frac{1}{T}(\int_{0}^{\frac{T}{2}}(f(t_0-\tau)-f(t_0^+))D_N(\tau)\,d\tau+\int_{-\frac{T}{2}}^{0}(f(t_0-\tau)-f(t_0^-))D_N(\tau)\,d\tau)         \\
                                           & =\frac{1}{T}\int_{T}g(t)(e^{i(N+1)\omega t}-e^{iN\omega t})\,dt                                                                                  \\
    \text{其中}g(t)                          & :=\begin{cases}
                                                   \frac{f(t_0+t)-f(t_0^-)}{e^{i\omega t}-1} & \text{if }-\frac{T}{2}<t_0<0 \\
                                                   \frac{f(t_0+t)-f(t_0^+)}{e^{i\omega t}-1} & \text{if }0<t_0<\frac{T}{2}
                                               \end{cases}                                                      \\
    \lim_{t\to 0^+}g(t)=                   & \lim_{t\to 0^+}\frac{f(t_0+t)-f(t_0^+)}{e^{i\omega t}-1}=\lim_{t\to 0^+}\frac{f(t_0+t)-f(t_0+)}{t}\lim_{t\to 0^+}\frac{t}{e^{i\omega t}-1}       \\
    |\lim_{t\to 0^+}g(t)|\leq              & |\frac{M}{i\omega}|=\frac{M}{\omega}
\end{align*}
$t\to 0^-$时同理有\[|\lim_{t\to 0^-}g(t)|\leq\frac{M}{\omega} \]
故g分段连续，当然是平方可积的，由命题（2.2），$g(t)$的傅里叶系数平方和收敛，通项趋于0，
\[S_N^f(t_0)-\frac{f(t_0^+)+f(t_0^-)}{2}=C_{-(N+1)}-C_N\to 0\]

\section{与狄利克雷条件的比较}
设函数$f(t)$周期为T。狄利克雷条件为：\begin{itemize}
    \item $\int_{T}|f(t)|\,dt<\infty$
    \item $f\in PC[0,T]$
    \item 在一个周期内，f的极值点个数有限
\end{itemize}
它是傅里叶级数存在且收敛至左右极限的平均值的充分条件。

根据命题（2.4），如果函数f满足\begin{itemize}
    \item $\int_{T}|f(t)|\,dt<\infty$
    \item $f\in PC[0,T]$
    \item 在点$t_0$处，四个迪尼导数有限
\end{itemize}
则可推出f的傅里叶级数存在且在$t_0$处收敛至左右极限的平均值。

满足狄利克雷条件的函数，迪尼导数未必有限：取周期为T的函数
\[f(t)=\begin{cases}
        \sqrt{t} & \text{if }0<t< \frac{T}{2}         \\
        0        & \text{if }-\frac{T}{2}\leq t\leq 0
    \end{cases}\]
则由于$\lim_{t\to 0^+}f(t)=0=f(0)$，f在0处连续，在一个周期内只有一个不连续点，
并且是此时这个不连续点是第一类间断点，f满足狄利克雷条件，但是
\begin{align*}
    D^+f(0)=\lim_{t\to 0^+}\frac{\sqrt{t}}{t}=\lim_{t\to 0^+}\frac{1}{\sqrt{t}}=\infty
\end{align*}

迪尼导数有限的函数未必满足狄利克雷条件：显然可以使函数满足狄利克雷条件的前两条，
并在一个周期内有无限多个极值点，因为迪尼导数有限是局部性质；即使在所有点迪尼导数
有限，也未必满足狄利克雷条件。例如取
\[f(t)=\begin{cases}
        t^2\sin\frac{1}{t} & \text{if }0<t<\frac{T}{2}          \\
        0                  & \text{if }-\frac{T}{2}\leq t\leq 0
    \end{cases}\]
则显然f连续，从而绝对可积，并且
\[\lim_{t\to 0^+}\frac{f(t)-f(0)}{t}=\lim_{t\to 0^+}t\sin\frac{1}{t}=0=\lim_{t\to 0^-}\frac{f(t)-f(0)}{t}\Rightarrow f'(0)=0\]
\[f'(t)=
    \begin{cases}
        2t\sin\frac{1}{t}-\cos\frac{1}{t} & \text{if }0<t<\frac{T}{2}      \\
        0                                 & \text{if }-\frac{T}{2}<x\leq 0
    \end{cases}\]
当$0<x<\frac{T}{2}$时，导数$f'(t)=2t\sin\frac{1}{t}-\cos\frac{1}{t}$
存在且有限，从而迪尼导数存在且有限。但是$t=\frac{1}{k\pi}(k\in\mathbb{N})$都
是f的零点，并且0的右邻域内零点均为此形式，区间$[\frac{1}{k\pi},\frac{1}{(k+1)\pi}]$
内函数不恒为0，必有极值点，所以f有无穷多个极值点，不满足狄利克雷条件。

综上，命题（2.4）中的条件与狄利克雷条件互不蕴含。

\end{document}
